% =====================================================================================
% REVISIÓN 1: solo los «Temas tentativos del marco teórico».
% REVISIÓN 2: el Marco teórico completo (cuando lo desarrolles, sustituye o convierte los temas tentativos en subsecciones).
% Escribe tu texto FUERA de los recuadros \begin{guia} ... \end{guia}.
% =====================================================================================

\section{Marco teórico} \label{sec:marco_teorico} % audit:req=marco_teorico

\begin{guia}[Guía: Marco teórico (Revisión 2)]
\textbf{¿Para qué sirve?} Reúne los conceptos, teorías, tecnologías, métodos y normas que se necesitan para entender y
justificar tu propuesta. No es un glosario ni un resumen de Internet: cada tema se explica \emph{en función de tu proyecto}.

\textbf{Debe responder (por cada tema):}
\begin{itemize}[nosep,leftmargin=*]
  \item ¿Qué es y cómo funciona? (definición con fuente citada y, si ayuda, un esquema o ecuación).
  \item ¿Para qué se usa en mi proyecto y por qué lo elegí frente a las alternativas?
  \item ¿Qué ventajas, limitaciones o buenas prácticas debo conocer?
\end{itemize}
Además: ¿qué metodología de desarrollo o investigación usarás y por qué? ¿qué normas o estándares aplican (calidad,
seguridad, protección de datos)?

\textbf{Organización:} una \verb|\subsection| por tema, de lo general a lo particular.

\textbf{Extensión mínima:} 400 palabras de desarrollo (sin contar los temas tentativos de la Revisión 1).

\textbf{Fuentes:} libros, artículos y documentación oficial. Cada afirmación importante lleva \verb|\cite{clave}|.

\textbf{Errores frecuentes:} temas sin relación con el proyecto; copiar y pegar; fuentes que son solo URLs de blogs; no
citar; temas que luego no se usan en el desarrollo.

\textbf{Cómo se revisa:} pertinencia, profundidad suficiente, calidad de las fuentes y que cada tema se vincule con la
propuesta.
\end{guia}

% >>> REVISIÓN 2: ESCRIBE AQUÍ el marco teórico (introducción breve y una \subsection por tema) <<<


\subsection{Temas tentativos del marco teórico} \label{sec:temas_tentativos} % audit:req=temas_marco_teorico

\begin{guia}[Guía: Temas tentativos del marco teórico (Revisión 1)]
\textbf{¿Para qué sirve?} En la primera revisión no se pide el marco teórico completo, sino la lista de los temas que
vas a desarrollar. Así el director puede confirmar que vas bien encaminado.

\textbf{Debe responder:} ¿cuáles son los temas (de 5 a 10) que necesito para sustentar el proyecto? Para cada uno, en una
o dos oraciones: ¿qué es y por qué lo necesito?

\textbf{Formato:} una lista (\verb|\begin{enumerate}| o \verb|itemize|), un tema por elemento.

\textbf{Extensión mínima:} 25 palabras. En la Revisión 2 conviertes cada tema en una subsección desarrollada.

\textbf{Cómo se revisa:} que los temas cubran lo necesario para el proyecto y no incluyan relleno.
\end{guia}

% >>> REVISIÓN 1: ESCRIBE AQUÍ la lista de temas tentativos <<<
\begin{enumerate}
    \item \textbf{Visión computacional:} emulación del reconocimiento visual humano mediante cómputo (detección, reconocimiento, seguimiento); se necesita para fundamentar el tratamiento de la imagen del documento como entrada del pipeline.
    \item \textbf{Procesamiento digital de imágenes:} conversión a escala de grises, binarización, corrección de inclinación (\textit{deskew}) y reducción de ruido; se necesita porque la calidad de la imagen de entrada condiciona directamente la legibilidad para el OCR.
    \item \textbf{Reconocimiento óptico de caracteres:} segmentación de líneas y palabras, reconocimiento de caracteres y confianza por palabra, incluyendo modos de segmentación (PSM) y el tratamiento de PDF con capa de texto frente a imagen escaneada; es el núcleo de la digitalización.
    \item \textbf{Clasificación de documentos:} clasificación por reglas y palabras clave (factura, orden\_compra, remision o desconocido) y sus métricas (matriz de confusión, precisión, cobertura); se necesita para tipificar cada documento antes de extraer sus campos.
    \item \textbf{Comprobantes fiscales en México (CFDI/SAT):} estructura del comprobante, RFC y UUID como identificadores, y campos fiscalmente relevantes (emisor, receptor, folio, fecha, conceptos, impuestos, total); se necesita porque la captura exacta de estos campos es el requisito fiscal del sistema.
    \item \textbf{Extracción estructurada y evaluación por campo:} identificación de campos con significado definido dentro de texto no estructurado, definición de acierto por campo, tolerancia de $\pm 1.00$ en montos y protocolo de evaluación sobre el conjunto mixto; se necesita para demostrar el desempeño del sistema.
\end{enumerate}

% >>> REVISIÓN 2: desarrollo del marco teórico (una \subsection por tema, de lo general a lo particular) <<<

Este marco teórico fundamenta cada etapa del pipeline implementado: la imagen del documento se preprocesa con OpenCV, se reconoce con Tesseract, se clasifica por reglas, se extrae a campos CFDI normalizados y se persiste mediante una API REST tipada y una base relacional, todo ello bajo un protocolo de métricas reproducible. Cada tema define el concepto, describe su funcionamiento técnico y lo vincula con la decisión adoptada en el sistema.

\subsection{Visión computacional aplicada a documentos} \label{sec:mt_vision}

La visión computacional es la disciplina que dota a la computadora de la capacidad de interpretar imágenes mediante modelos matemáticos de formación de imagen, extracción de características y reconocimiento de patrones\cite{bradski_opencv_2008}. Su objeto no es reproducir la visión humana, sino construir representaciones operables de la escena —bordes, regiones, líneas de texto— sobre las cuales decidir con algoritmos deterministas o aprendidos\cite{bradski_opencv_2008, opencv_docs}.

En el presente proyecto, la visión computacional se restringe al dominio documental: la entrada no es una escena arbitraria sino una página con estructura de layout (encabezado, tablas de conceptos, pie con totales). Esta restricción es una ventaja metodológica, pues permite suponer orientación predominantemente horizontal del texto, fondo aproximadamente uniforme y bloques de lectura ordenados, supuestos que el preprocesamiento y el modo de segmentación del OCR explotan directamente\cite{smith_tesseract_2007, opencv_docs}. Frente a la alternativa de tratar cada documento como imagen genérica con redes profundas de propósito general, el enfoque documental clásico ofrece auditabilidad etapa por etapa y ejecución en CPU local sin aceleradores, coherente con los recursos del proyecto\cite{bradski_opencv_2008, tesseract_ocr}.

\subsection{Procesamiento digital de imágenes con OpenCV} \label{sec:mt_pdi}

El procesamiento digital de imágenes transforma la imagen captada en una representación con mayor legibilidad para las etapas posteriores, sin alterar su contenido semántico\cite{bradski_opencv_2008, opencv_docs}. El pipeline adoptado consta de tres operaciones encadenadas. Primero, la conversión a escala de grises reduce cada píxel a una luminancia única, lo cual elimina la variabilidad cromática de logotipos y sellos y concentra la información en el contraste trazo-fondo\cite{bradski_opencv_2008}. Segundo, la corrección de inclinación (\textit{deskew}) estima el ángulo de rotación de las líneas de texto —típicamente mediante transformada de Hough o análisis de momentos— y aplica una transformación afín inversa que horizontaliza los renglones, condición que los segmentadores de línea del OCR presuponen\cite{bradski_opencv_2008, opencv_docs}. Tercero, la binarización adaptativa calcula un umbral local por vecindad en lugar de un umbral global, de modo que sombras, fondos degradados y densidades heterogéneas de tóner se separan correctamente en texto y fondo\cite{bradski_opencv_2008, opencv_docs}.

En el sistema, estas operaciones se implementan con OpenCV porque la biblioteca documenta explícitamente los módulos de procesamiento (conversión de color, umbralización y transformaciones geométricas) y porque su ejecución en CPU satisface los tiempos del pipeline sin dependencias propietarias\cite{opencv_docs}. La limitación conocida es que un preprocesamiento agresivo puede fragmentar trazos finos o unir caracteres adyacentes; por ello los parámetros se calibran sobre el conjunto de documentos del proyecto y no se asumen universales\cite{bradski_opencv_2008, opencv_docs}.

\subsection{Reconocimiento óptico de caracteres con Tesseract} \label{sec:mt_ocr}

El reconocimiento óptico de caracteres (OCR) convierte imágenes de texto en caracteres codificados mediante una secuencia de segmentación de líneas y palabras, reconocimiento de cada símbolo con modelos entrenados y asignación de una confianza a cada decisión\cite{smith_tesseract_2007}. Tesseract, el motor abierto adoptado, expone dos controles determinantes para este trabajo: el idioma de reconocimiento y el modo de segmentación de página (PSM). El sistema opera con idioma español (\textit{spa}) y PSM~6, modo que presupone un bloque uniforme de texto y delega la segmentación interna de líneas, supuesto válido para páginas de factura ya enderezadas\cite{tesseract_ocr, smith_tesseract_2007}.

Cada palabra reconocida lleva asociada una confianza (\textit{ocr\_confidence}) que el sistema propaga hasta el esquema de datos; los campos con confianza baja se marcan para revisión humana dirigida en lugar de obligar a revisar el documento completo\cite{tesseract_ocr}. Adicionalmente, el pipeline distingue dos naturalezas de PDF: los documentos con capa de texto se leen directamente con extracción de texto (pypdf), sin pasar por visión, mientras que los escaneados sin capa se renderizan a imagen y siguen la ruta OpenCV más Tesseract\cite{tesseract_ocr, opencv_docs}. Frente a servicios de OCR en la nube, que ofrecen mayor exactitud base a cambio de costo por documento y de enviar comprobantes fiscales a infraestructura externa, Tesseract local preserva la confidencialidad de los documentos, elimina el costo marginal y permite auditar cada etapa, a costa de exigir preprocesamiento cuidadoso\cite{gcp_document_ai, tesseract_ocr}.

\subsection{Clasificación de documentos por reglas} \label{sec:mt_clasificacion}

La clasificación de documentos asigna cada texto reconocido a una categoría del vocabulario operativo del sistema: \textit{factura}, \textit{orden\_compra}, \textit{remision} o \textit{desconocido}\cite{sokolova_metrics_2009}. El clasificador implementado es determinista y lexical: aplica expresiones regulares y listas de palabras clave propias de cada tipo documental (por ejemplo, marcadores de folio fiscal frente a marcadores de orden o entrega) y resuelve empates y ausencias hacia la clase \textit{desconocido}, de modo que ningún documento queda silenciosamente mal tipificado\cite{sokolova_metrics_2009, structured_extraction_2023}.

La elección de reglas frente a un clasificador aprendido se justifica por tres razones vinculadas al contexto. Primero, el vocabulario discriminante de estos documentos es explícito y estable (las facturas mexicanas exhiben marcas fiscales normalizadas), por lo que las reglas capturan la señal sin necesidad de corpus de entrenamiento\cite{sat_cfdi_2024}. Segundo, las reglas son inspeccionables: cada decisión puede explicarse por la regla disparada, propiedad exigible en un sistema con efectos fiscales\cite{sat_cfdi_validation_2023}. Tercero, su costo computacional es despreciable frente al OCR, como lo confirma la medición del proyecto sobre texto limpio\cite{sokolova_metrics_2009}. Su limitación —fragilidad ante formatos con vocabulario inédito— se mitiga con la clase \textit{desconocido} y con la evaluación por clase descrita en la Sección~\ref{sec:mt_metricas}\cite{structured_extraction_2023}.

\subsection{Comprobante Fiscal Digital por Internet (CFDI) y SAT} \label{sec:mt_cfdi}

El Comprobante Fiscal Digital por Internet (CFDI) es el esquema de facturación electrónica administrado por el Servicio de Administración Tributaria (SAT): cada comprobante es un documento XML timbrado conforme al Anexo~20, identificado por un folio fiscal único (UUID) y con emisor y receptor identificados por su Registro Federal de Contribuyentes (RFC)\cite{sat_cfdi_2024}. Los campos fiscalmente relevantes para este trabajo son el RFC del emisor, el RFC del receptor, el folio, la fecha de emisión, los conceptos (renglones de la operación), los impuestos trasladados y el total; la validez de la operación depende de la coincidencia exacta de estos identificadores con lo timbrado\cite{sat_cfdi_2024, sat_cfdi_validation_2023}.

El sistema trata RFC y UUID como ciudadanos de primera clase del esquema de datos: poseen validadores de formato propios y metas de exactitud diferenciadas, porque un error de un solo carácter en ellos invalida la trazabilidad fiscal del comprobante mientras que un error equivalente en un campo descriptivo no\cite{sat_cfdi_2024}. El SAT ofrece además el servicio de verificación de comprobantes, que constituye la referencia externa contra la cual la organización puede contrastar los campos extraídos\cite{sat_cfdi_validation_2023}. Esta especificidad fiscal mexicana es precisamente lo que distingue al sistema de extractores genéricos internacionales, que no modelan el UUID ni el RFC como campos nativos\cite{gcp_document_ai, sat_cfdi_2024}.

\subsection{Extracción estructurada y normalización} \label{sec:mt_extraccion}

La extracción estructurada identifica, dentro de texto no estructurado, un conjunto de campos con significado definido (folio, fecha, RFC, UUID, conceptos, totales) y los representa conforme a un esquema tipado y validable\cite{structured_extraction_2023}. En documentos de compra, la tarea se complica porque los formatos varían entre emisores, los conceptos aparecen como tablas de renglones múltiples y los montos admiten separadores y símbolos heterogéneos; por ello el sistema combina expresiones regulares por campo con heurísticas de contexto (vecindad de etiquetas como ``total'' o ``folio'') y encapsula el resultado en un objeto \textit{ExtractionResult} que incluye advertencias de integridad, entre ellas la de descuadre de totales (\textit{TOTALS\_MISMATCH}) cuando subtotal más impuestos difiere del total declarado\cite{structured_extraction_2023, kim_donut_2022}.

La normalización posterior convierte cada valor extraído a su forma canónica: fechas a ISO~8601, montos a decimal con moneda MXN explícita y cadenas fiscales a mayúsculas sin espacios espurios, de modo que la comparación y la persistencia sean deterministas\cite{structured_extraction_2023, pydantic_v2_2024}. Frente a la alternativa de delegar la extracción a un modelo multimodal sin OCR, que promete robustez ante layouts inéditos a cambio de cómputo intensivo y comportamiento menos auditable, las reglas normalizadas ofrecen trazabilidad total campo por campo y constituyen la base medible sobre la cual una fase futura con modelo de lenguaje podrá compararse\cite{kim_donut_2022, structured_extraction_2023}.

\subsection{Backend: FastAPI, Pydantic v2, PostgreSQL y SQLAlchemy} \label{sec:mt_backend}

El backend expone el pipeline mediante una API REST, estilo arquitectónico que organiza el contrato alrededor de recursos (documentos, resultados) y operaciones estándar de HTTP, con separación entre cliente y servidor\cite{fielding_rest_2000}. FastAPI implementa dicho contrato con enrutadores tipados, validación automática de solicitudes y documentación interactiva generada a partir de los esquemas, lo cual reduce los errores de integración con los sistemas consumidores\cite{fastapi_docs}. La validación estricta se delega a Pydantic~v2, que verifica tipos, formatos (RFC, UUID, ISO~8601) y restricciones de rango antes de que cualquier dato alcance la base, de modo que el contrato de la API es ejecutable y no meramente documental\cite{pydantic_v2_2024, fastapi_docs}.

La persistencia se implementa sobre PostgreSQL~16, gestor relacional que garantiza atomicidad, consistencia, aislamiento y durabilidad (ACID) mediante control de concurrencia multiversión, tipos de datos estrictos e índices que sostienen la consulta por folio, RFC y UUID\cite{postgres_docs_2024}. SQLAlchemy actúa como capa de mapeo objeto-relacional: los modelos declarativos (\textit{Document}, páginas, campos, resultados) definen el esquema en Python tipado y las sesiones transaccionales aseguran que cada documento y sus derivados se persistan de forma atómica\cite{sqlalchemy_docs_2024}. El procesamiento pesado (OCR, clasificación, extracción) se ejecuta en segundo plano con tareas asíncronas (\textit{BackgroundTasks}), con estados \textit{processing}, \textit{done} y \textit{error}, de modo que la carga del archivo responde de inmediato y el cliente consulta el resultado cuando el worker lo concluye\cite{fastapi_docs, postgres_docs_2024}.

\subsection{Contenerización con Docker} \label{sec:mt_docker}

La contenerización empaqueta la aplicación con sus dependencias y configuración en imágenes reproducibles que se ejecutan como contenedores aislados, orquestados declarativamente con Docker Compose\cite{docker_docs}. En este proyecto, Compose define dos servicios —\textit{api} y \textit{db}— con verificaciones de salud (\textit{healthchecks}): la base de datos declara su disponibilidad y la API solo inicia cuando la base responde, además de exponer su propio chequeo de vida; los volúmenes persistentes conservan los datos de PostgreSQL y los archivos cargados entre reinicios\cite{docker_docs}.

Frente a la alternativa de despliegue directo sobre el sistema operativo (\textit{bare-metal}), la contenerización garantiza que el entorno evaluado (versión de Tesseract con idioma español, OpenCV, PostgreSQL~16) sea idéntico en cualquier máquina, condición necesaria para que las métricas reportadas sean reproducibles\cite{docker_docs}. La limitación es el sobrecosto de virtualización ligera y la necesidad de construir la imagen con los datos de idioma del OCR incluidos; por ello la medición de OCR real sobre documentos escaneados queda explícitamente condicionada al entorno Docker completo, como se detalla en la propuesta\cite{docker_docs, tesseract_ocr}.

\subsection{Métricas de exactitud y latencia} \label{sec:mt_metricas}

La evaluación del sistema es por campo y no por documento: cada campo extraído (RFC, UUID, folio, fecha, cada concepto, cada monto) se compara contra su valor de referencia y solo el acierto exacto cuenta, salvo en montos donde se admite una tolerancia de $\pm 1.00$ por redondeos de representación\cite{sokolova_metrics_2009, structured_extraction_2023}. Sobre esta definición se construyen las métricas de clasificación: la precisión (fracción de extracciones correctas entre las emitidas), la cobertura (fracción de campos de referencia recuperados) y la medida F como media armónica de ambas, que penaliza los sistemas que sacrifican una por la otra\cite{sokolova_metrics_2009}. Dado que las clases documentales están desbalanceadas (predominan las facturas), el promedio macro —media simple de la medida F por clase— evita que la clase mayoritaria oculte el desempeño en órdenes de compra y remisiones\cite{sokolova_metrics_2009}.

Dos métricas adicionales capturan requisitos no funcionales del contrato: la tasa de JSON válido (fracción de respuestas del pipeline que satisfacen el esquema \textit{ExtractionResult} sin violaciones de Pydantic) y los percentiles de latencia p50 y p95 de la etapa de reglas, que caracterizan el tiempo típico y el de cola respectivamente\cite{pydantic_v2_2024, sokolova_metrics_2009}. El protocolo se ejecuta sobre un conjunto de 30 documentos oficiales más 4 excluidos documentados, con scripts de evaluación y de medición (\textit{bench}) versionados; al redactar este avance, la etapa de reglas alcanza 100\% sobre texto limpio con p50 de 0.08~ms y p95 de 0.09~ms, mientras que la medición con OCR real queda pendiente del entorno Docker completo\cite{structured_extraction_2023}.

% EVIDENCIA: bradski_opencv_2008 --- anexar portada e índice del libro Learning OpenCV (O'Reilly, 2008) que respalden conversión de color, umbralización y transformaciones afines.
% EVIDENCIA: opencv_docs --- anexar impresión de las páginas de imgproc (threshold, moments, transformaciones geométricas) con fecha de consulta.
% EVIDENCIA: smith_tesseract_2007 --- anexar portada del acta ICDAR 2007 y resumen del motor Tesseract (segmentación, reconocimiento, confianza).
% EVIDENCIA: tesseract_ocr --- anexar impresión del repositorio oficial y de la wiki (modos PSM, idioma spa, confianza por palabra) con fecha de consulta.
% EVIDENCIA: sat_cfdi_2024 --- anexar impresión de la página oficial del CFDI y del Anexo 20 vigente con fecha de consulta.
% EVIDENCIA: sat_cfdi_validation_2023 --- anexar impresión del portal de verificación de CFDI del SAT con fecha de consulta.
% EVIDENCIA: fastapi_docs --- anexar impresión de la documentación de FastAPI (BackgroundTasks, validación de esquemas) con fecha de consulta.
% EVIDENCIA: pydantic_v2_2024 --- anexar impresión de la documentación de Pydantic v2 (validadores estrictos) con fecha de consulta.
% EVIDENCIA: postgres_docs_2024 --- anexar impresión de la documentación de PostgreSQL 16 (ACID, MVCC, índices) con fecha de consulta.
% EVIDENCIA: sqlalchemy_docs_2024 --- anexar impresión de la documentación de SQLAlchemy (mapeo declarativo, sesiones) con fecha de consulta.
% EVIDENCIA: docker_docs --- anexar impresión de la documentación de Docker y Compose (imágenes, healthchecks, volúmenes) con fecha de consulta.
% EVIDENCIA: fielding_rest_2000 --- anexar portada e índice de la tesis doctoral de Fielding (REST, UC Irvine, 2000).
% EVIDENCIA: sokolova_metrics_2009 --- anexar portada del artículo (Inf. Process. Manag. 45(4), DOI 10.1016/j.ipm.2009.03.002) con definiciones de precisión, cobertura, F y promedios.
% EVIDENCIA: structured_extraction_2023 --- anexar portada del pre-print LayoutLM (arXiv 1912.13318) que respalda extracción con layout.
% EVIDENCIA: kim_donut_2022 --- anexar portada del acta ECCV 2022 / pre-print (arXiv 2111.15664) que respalda la fase futura sin OCR.
% EVIDENCIA: gcp_document_ai --- anexar impresión de la documentación de Document AI (facturas, tablas, precios por documento) con fecha de consulta.
